\documentclass[10pt,a4paper]{amsart}
\usepackage[margin=19mm]{geometry}
\usepackage{amsmath,amssymb,mathtools}
\usepackage[scr=rsfs,cal=euler]{mathalfa}
\usepackage{xfrac}
\usepackage[unicode,hidelinks,bookmarks=false]{hyperref}
\newcommand{\ie}{i.e.\ }
\newcommand{\cf}{cf.\ }
\newcommand{\resp}{respectively\ }
\newcommand{\Subsection}{\subsection}
\providecommand{\nobreakdash}{\nobreak\hbox{-}}
\setlength{\emergencystretch}{2em}
\multlinegap0pt
\usepackage{amssymb}
\usepackage{xcolor}
\usepackage[normalem]{ulem}
\usepackage{tikz}
\newtheorem{lemma}{Lemma}
\title{\small Rings whose non-invertible elements are \\ uniquely strongly clean}
\author{Peter Danchev, Mehrdad Esfandiar, Omid Hasanzadeh}
\date{Source: arXiv:2401.03320v2 (CC BY 4.0)}
\begin{document}
\maketitle

\noindent\emph{Source and licence.} \small Rings whose non-invertible elements are \\ uniquely strongly clean. Version arXiv:2401.03320v2 (CC BY 4.0), distributed by arXiv under the Creative Commons Attribution 4.0 International licence, http://creativecommons.org/licenses/by/4.0/, as recorded in the arXiv licence metadata for this exact version and shown on the abstract page https://arxiv.org/abs/2401.03320v2. Full licence notice: \texttt{licenses/arxiv-2401.03320-CC-BY-4.0.txt}.\par\smallskip

\noindent\textit{Editorial scope.} Counted unit: the article's lemma 2.6 (source0.tex lines 275--277). The article's complete original proof of it is delivered with it (source0.tex lines 279--297, one \texttt{proof} environment of 19 lines). No mathematics is written, repaired or re-derived: the statement and the proof are the article's own lines, in the article's own notation, and no retained line is substituted or re-flowed.  No further source-local context is needed: every cross-reference of every retained line resolves inside the two delivered spans.  Nothing was omitted from any retained span, so the package carries no omission record.  No retained line contains a citation, so the wrapper carries no bibliography and no bibliography entry.  No retained line contains a figure environment, an image-inclusion command or drawing code, so no image is delivered and the package declares no asset dependency.  Editorial formatting only, in addition: an AMS \texttt{amsart} preamble that restates the article's own declarations and macros used by the retained lines, namely the environments \texttt{lemma} on separate counters, together with the standard packages those lines need.  The retained environments print in this document as Lemma 1 in order of appearance, whereas the article prints the same items with the numbering of its own full document.  The complete native source of the article as distributed in the arXiv e-print for this exact version is bundled unchanged as \texttt{sources/152567/source0.tex}.
\begin{lemma}\label{lemma 2.6}
Let $R$ be a GUSC ring with $2 \in U(R)$ and, for every $u \in U(R)$, we have $u^2 = 1$. Then, $R$ is a commutative ring.
\end{lemma}

\begin{proof}
Firstly, we show that \( R \) is reduced, i.e., \( R \) contains no non-trivial nilpotent elements, so suppose \( q \in Nil(R) \). Then, \( (1 \pm q) \in U(R) \), whence
\[ 1 - 2q + q^2 = (1-q)^2 = 1 = (1+q)^2 = 1 + 2q + q^2. \]
Thus, \( 4q = 0 \). Since \( 2 \in U(R) \), we conclude \( q = 0 \). Hence, \( R \) is reduced and, consequently, \( R \) is abelian.
	
Moreover, for any \( u, v \in U(R) \), we have \( u^2 = v^2 = (uv)^2 = 1 \). Therefore, \( uv = (uv)^{-1} = v^{-1} u^{-1} = vu \), and hence the units commute with each other.
	
In addition, let \( x, y \in R \). We consider the following four possible cases:
	
1. \( x, y \in U(R) \): since the units commute, it must be that \( xy = yx \).
	
2. \( x, y \notin U(R) \): since \( R \) is a GUSC ring, there exist \( e, f \in Id(R) \) and $u, v\in U(R)$ such that \( x = e+u \) and \( y = f+v \). So, \( xy = yx \), because \( R \) is abelian and the units commute.
	
3. \( x \in U(R) \) and \( y \notin U(R) \): in this case, there exist \( e \in Id(R) \) and $u\in U(R)$ such that \( y = e+u \). Then, $xy=x(e+u)=xe+xu$. But, we have $xy=ex+ux=(e+u)x=yx$, because $R$ is abelian and the units commute.
	
4. \( x \notin U(R) \) and \( y \in U(R) \): similarly to case (3), we can easily see that \( xy = yx \).	

This completes our arguments after all.
\end{proof}

\end{document}
